\section{One Marian confession, several kinds of claim}

Catholic speech about Mary ranges across solemn dogma, authoritative doctrine,
liturgical confession, patristic typology, received devotion, theological
opinion, historical report, and private revelation. These are not rival
Maries. They are different acts of knowing, teaching, remembering, and praying
within the one mystery of Christ. Confusing them harms both faith and devotion:
a dogma becomes a slogan, a title becomes a metaphysical mechanism, an
apparition becomes a fifth Gospel, or a legitimate opinion becomes a test of
Catholic fidelity.

The Church's distinctions serve natural goods: clarity about who claims what,
protection against equivocation, historical memory, fair disagreement, and a
conscience not burdened by invented obligations. Their supernatural end is to
keep Marian doctrine ordered to the Incarnation and redemption, foster the
obedience of faith, protect Christ's unique mediation, and guide devotion into
Scripture, liturgy, sacrament, conversion, charity, and ecclesial communion.

